% !TEX root = ../Presentation.tex
% ---------------------------------------------------------------------------
% Model B: hierarchical over sessions
% Slides taken from ../Presentation_ABC/sections; trimmed ones are marked.
% ---------------------------------------------------------------------------
\section{Model B: Sessions Hierarchical}

\subsection{Model Structure}

\begin{frame}{Model B: graphical model}
  \centering
  \begin{tikzpicture}[gm, scale=0.82, transform shape]
    \foreach \n/\lab [count=\k from 0] in {mub0/\mu_{\beta_0}, sb0/\sigma_{\beta_0},
        mub1/\mu_{\beta_1}, sb1/\sigma_{\beta_1}, mugh/\mu_{\gamma_h}, sgh/\sigma_{\gamma_h},
        mugl/\mu_{\gamma_l}, sgl/\sigma_{\gamma_l}}
      \node[latent] (\n) at (1.1*\k, 0) {$\lab$};

    \node[latent] (b0) at ($(mub0)!0.5!(sb0) + (0,-1.6)$) {$\beta_{0,s}$};
    \node[latent] (b1) at ($(mub1)!0.5!(sb1) + (0,-1.6)$) {$\beta_{1,s}$};
    \node[latent] (gh) at ($(mugh)!0.5!(sgh) + (0,-1.6)$) {$\gamma_{h,s}$};
    \node[latent] (gl) at ($(mugl)!0.5!(sgl) + (0,-1.6)$) {$\gamma_{l,s}$};

    \node[det] (psi) at ($(b1)!0.5!(gh) + (0,-1.4)$) {$\psi_i$};
    \node[obs, below=0.5cm of psi] (r) {$r_i$};
    \node[const, right=0.8cm of psi] (x) {$x_i$};

    \foreach \h/\p in {mub0/b0, sb0/b0, mub1/b1, sb1/b1, mugh/gh, sgh/gh, mugl/gl, sgl/gl}
      \draw[arr] (\h) -- (\p);
    \foreach \p in {b0, b1, gh, gl}
      \draw[arr] (\p) -- (psi);
    \draw[arr] (x) -- (psi);
    \draw[arr] (psi) -- (r);

    \node[plate, fit=(x)(psi)(r), inner sep=5pt] (ptrial) {};
    \node[font=\scriptsize, anchor=north east] at (ptrial.south east) {trials $i$ in session $s$};
    \node[plate, fit=(b0)(gl)(ptrial), inner sep=12pt] (psess) {};
    \node[font=\scriptsize, anchor=south west] at (psess.south west) {sessions $s = 1,\dots,42$};
    \node[plate, fit=(mub0)(sgl)(psess), inner sep=8pt] (pgroup) {};
    \node[font=\scriptsize, anchor=north east] at (pgroup.south east) {each of the 4 groups};
  \end{tikzpicture}

  \vfill
  {\scriptsize\color{muted} Circles: random variables (shaded: observed); rectangle:
  deterministic; $x_i$: fixed covariate. The same model is fit independently to each group.}
\end{frame}

\subsection{Synthetic Data}
% numbers from Model_B/Convergence_B_Report.tex

% trimmed
\begin{frame}{Model B on synthetic data: what is recoverable}
  At the size of the real data (42 sessions of about 308 trials), at values like the real ones:
  \begin{itemize}
    \item \textbf{Recoverable:} each session's PSE and JND, and how much they vary between
          sessions
    \item \textbf{Weakly determined:} the lapse rates, especially session by session
    \item \textbf{Not identified:} how much the lapse rates vary between sessions; the
          estimate reflects the prior
  \end{itemize}
\end{frame}

\subsection{Real Data}
% numbers from Model_B/Fit_B_Report.tex

\begin{frame}{Model B on the real data: fitted curves in one session (04-13)}
  \centering
  \includegraphics[width=0.48\textwidth]{\plotsB/Fit_Plots_B/curves_left_04-13.png}\hfill
  \includegraphics[width=0.48\textwidth]{\plotsB/Fit_Plots_B/curves_right_04-13.png}
\end{frame}

\begin{frame}{Model B on the real data: PSE, left hand}
  \centering
  \begin{tikzpicture}
      \node[inner sep=0, anchor=north west] (ub) at (-17.87bp,0)
        {\includegraphics[scale=0.3813, viewport=0 266.04 1016.28 487.08, clip]{\plotsB/Fit_Plots_B/PSE_uni_bi_by_session.png}};
      \node[inner sep=0, anchor=north west] (ef) at (-28.05bp,-87.60bp)
        {\includegraphics[scale=0.3896, viewport=0 266.04 1016.28 487.08, clip]{\plotsB/Fit_Plots_B/effect_PSE_by_session.png}};
      \node[inner sep=0, anchor=north west] (tk) at (-28.05bp,-173.70bp)
        {\includegraphics[scale=0.3896, viewport=0 0.00 1016.28 45.72, clip]{\plotsB/Fit_Plots_B/effect_PSE_by_session.png}};
  \end{tikzpicture}
\end{frame}

\begin{frame}{Model B on the real data: PSE, right hand}
  \centering
  \begin{tikzpicture}
      \node[inner sep=0, anchor=north west] (ub) at (-17.87bp,0)
        {\includegraphics[scale=0.3813, viewport=0 45.72 1016.28 266.04, clip]{\plotsB/Fit_Plots_B/PSE_uni_bi_by_session.png}};
      \node[inner sep=0, anchor=north west] (ef) at (-28.05bp,-87.32bp)
        {\includegraphics[scale=0.3896, viewport=0 45.72 1016.28 266.04, clip]{\plotsB/Fit_Plots_B/effect_PSE_by_session.png}};
      \node[inner sep=0, anchor=north west] (tk) at (-28.05bp,-173.14bp)
        {\includegraphics[scale=0.3896, viewport=0 0.00 1016.28 45.72, clip]{\plotsB/Fit_Plots_B/effect_PSE_by_session.png}};
  \end{tikzpicture}
  \par\vspace{0.1em}{\tiny\color{muted} \textcolor{plotblue}{$\bullet$}~unimanual, \textcolor{plotred}{$\bullet$}~bimanual: session mean and 95\% HDI; bands: typical session; green dashed: trained threshold\\ bi $-$ uni: filled if the session's 95\% HDI excludes 0; orange band: average over sessions}
\end{frame}

\plotframe{Model B: PSE shift vs distractor amplitude (exploratory)}
  {\plotsB/Fit_Plots_B/effect_vs_distractor_amplitude.png}

% trimmed
\begin{frame}{Model B: conclusions}
  \begin{itemize}
    \item Reliable fit; no systematic misfit, pooled or within sessions
    \item Effect of the distractor (bimanual $-$ unimanual), average over sessions:
    \begin{itemize}
      \item \textbf{PSE:} lowered, by about 4~$\mu m$ (left) and 6~$\mu m$ (right), in nearly
            every session
      \item \textbf{JND:} raised for the left hand, driven by a few sessions (possibly lapses
            absorbed into the slope)
      \item \textbf{\boldmath$\gamma_h$:} raised; \textbf{\boldmath$\gamma_l$:} no effect
    \end{itemize}
    \item The PSE shift grows with distractor amplitude (exploratory; amplitude is
          confounded with session order)
    \item[$\Rightarrow$] Put the distractor amplitude into the model: Model C
  \end{itemize}
\end{frame}
